%%
%% JPDC 2026 submission draft — v2
%% Sections 3-4 complete; other sections placeholder
%%
\documentclass[a4paper,fleqn]{cas-dc}
\usepackage[authoryear,longnamesfirst]{natbib}
\usepackage{amsmath,amssymb,amsfonts}
\usepackage{algorithm,algorithmic}
\usepackage{graphicx}
\usepackage{booktabs}
\usepackage{subcaption}
\usepackage{xcolor}
\usepackage{url}

% Theorem environments (separate counters so Thm starts at 1)
\newtheorem{theorem}{Theorem}
\newtheorem{proposition}{Proposition}
\newtheorem{corollary}{Corollary}
\newdefinition{definition}{Definition}
\newproof{pf}{Proof}

% Notation shortcuts
\newcommand{\ES}{\mathrm{ES}}
\newcommand{\Dref}{\mathcal{D}_{\mathrm{ref}}}
\newcommand{\Dpriv}{\mathcal{D}_i^{\mathrm{priv}}}
\newcommand{\Mset}{\mathcal{M}}
\newcommand{\Sset}{\mathcal{S}}
\newcommand{\Srecv}{\mathcal{S}^{\mathrm{recv}}}
\newcommand{\vmax}{v_{\max}}
\newcommand{\Bres}{B_{\mathrm{res}}}

\begin{document}
	\let\WriteBookmarks\relax
	\def\floatpagepagefraction{1}
	\def\textpagefraction{.001}
	
	% ── Metadata ──────────────────────────────────────────────
	\shorttitle{DASH for Asynchronous Edge Intelligence}
	\shortauthors{[Authors]}
	
	\title[mode=title]{Straggler-Tolerant Knowledge Distillation
		Scheduling for Asynchronous Edge Intelligence Systems}
	
	\author[1]{[First Author]}
	\cormark[1]
	\ead{[email]}
	\affiliation[1]{organization={Department of Computer Science
			and Information Engineering, Tamkang University},
		city={New Taipei City},
		postcode={251301},
		country={Taiwan}}
	\cortext[1]{Corresponding author}
	
	% ── Abstract ──────────────────────────────────────────────
	\begin{abstract}
		Federated distillation (FD) enables collaborative edge learning by
		exchanging soft predictions rather than model parameters, offering
		communication efficiency and architectural flexibility.  However,
		deploying FD over heterogeneous wireless edge networks faces a
		critical yet overlooked challenge: device stragglers whose
		unpredictable latencies dominate synchronous round times and degrade
		system throughput.  Existing asynchronous federated learning methods
		target parameter-level aggregation and cannot be directly applied to
		logit-based distillation, where partial uploads and quality-aware
		scheduling introduce unique design constraints.  This paper presents
		DASH (Deadline-Aware Straggler-Handling), an asynchronous
		distillation scheduling framework with straggler tolerance.  We
		design a three-outcome straggler model (complete, partial, timeout)
		driven by log-normal latency distributions, and integrate device
		completion probabilities into a submodular greedy selection
		algorithm.  We prove that the resulting expected quality retains a
		$(1-\bar{p})\tfrac{1}{2}(1-1/e)$ approximation ratio relative to the
		deterministic optimum, where $\bar{p}$ is the worst-case straggler
		probability.  Three timeout policies (fixed, adaptive, and
		partial-accept) provide system-level trade-offs between round latency
		and knowledge utilization.  Experiments on CIFAR-100 and EMNIST with
		up to 200 devices demonstrate that DASH achieves comparable accuracy
		to synchronous scheduling while reducing wall-clock convergence time
		by $2$--$3\times$ under 30\% straggler conditions, with graceful
		degradation as straggler severity increases.
	\end{abstract}
	
	% ── Highlights ────────────────────────────────────────────
	\begin{highlights}
		\item Asynchronous federated distillation protocol with a
		three-outcome straggler model for edge networks
		\item Straggler-aware submodular device selection with provable
		expected quality bound
		\item Three timeout policies (fixed, adaptive, partial-accept) with
		system-level trade-off analysis
		\item Large-scale experiments on two datasets with up to 200 devices
		and wall-clock time evaluation
	\end{highlights}
	
	% ── Keywords ──────────────────────────────────────────────
	\begin{keywords}
		Federated distillation \sep Asynchronous scheduling \sep
		Edge intelligence \sep Straggler tolerance \sep
		Submodular optimization
	\end{keywords}
	
	\maketitle
	
	% ══════════════════════════════════════════════════════════
	\section{Introduction}\label{sec:intro}
	% ══════════════════════════════════════════════════════════
	
	% [PLACEHOLDER — to be written after experiments]
	
	% ══════════════════════════════════════════════════════════
	\section{Related Work}\label{sec:related}
	% ══════════════════════════════════════════════════════════
	
	% [PLACEHOLDER — 2.1 Federated Distillation, 2.2 Async FL,
	%  2.3 Client Selection and Straggler Mitigation]
	
	% ══════════════════════════════════════════════════════════
	\section{System Architecture}\label{sec:system}
	% ══════════════════════════════════════════════════════════
	
	We consider an edge intelligence platform in which a single edge
	server~(ES) coordinates federated distillation across a pool
	$\Mset = \{1, \dots, M\}$ of heterogeneous IoT devices connected
	through unreliable wireless links.  Rather than exchanging model
	parameters, the ES operates a knowledge distillation service pipeline:
	in each round, a selected subset of devices compute soft predictions
	(logits) on a shared reference dataset and upload them to the ES,
	which distills the collected knowledge into a global model.  The
	central challenge addressed in this paper is that device latencies are
	heterogeneous and stochastic, so that synchronous protocols---which
	wait for all selected devices---suffer from worst-case round times
	that grow with the pool size~$M$.
	
	% ── 3.1 System Overview ──────────────────────────────────
	\subsection{System overview}\label{sec:overview}
	
	\begin{figure}[t]
		\centering
		% TODO: replace with actual figure file
		\fbox{\parbox{0.92\linewidth}{\centering\vspace{3.5cm}%
				\textit{Fig.~1: System architecture diagram}\\[2pt]
				\small ES modules (DASH Scheduler $\to$ Async Controller
				$\to$ Aggregation Engine) connected to $M$ devices
				via unreliable channels.  Three device outcomes shown:
				on-time~(\checkmark), partial~($\triangle$),
				timeout~($\times$).
				\vspace{0.5cm}}}
		\caption{DASH system architecture.  The ES comprises three
			modules: the DASH Scheduler (volume allocation and
			straggler-aware device selection), the Async Controller
			(deadline enforcement and logit collection), and the
			Aggregation Engine (quality-weighted distillation).  Each
			device may finish on time, partially upload, or timeout.}
		\label{fig:system}
	\end{figure}
	
	The DASH architecture comprises three server-side modules and a
	device-side execution pipeline, as illustrated in
	Fig.~\ref{fig:system}.
	
	On the server side, the DASH Scheduler determines at the beginning of
	each round which devices to engage and how many logit vectors to
	request from each, subject to a per-round resource budget~$B$.  The
	Async Controller sets a per-round deadline~$D^{(t)}$, dispatches
	upload quotas to selected devices, and collects responses as they
	arrive; unlike synchronous protocols that block until all devices
	respond, the Async Controller closes the collection window at
	$D^{(t)}$ regardless of how many devices have reported.  The
	Aggregation Engine receives the collected logits, computes
	quality-weighted averages on the reference dataset~$\Dref$, and
	updates the global model through knowledge distillation.
	
	On the device side, upon receiving the global model
	$\mathbf{w}^{(t-1)}$ and its upload quota~$v_i^*$, each selected
	device~$i$ fine-tunes its local model on its private
	dataset~$\Dpriv$, computes logits on~$\Dref$ and applies local
	perturbation, and uploads up to~$v_i^*$ perturbed logit vectors to
	the ES\@.  The total time for this pipeline depends on the device's
	computational capability, channel conditions, and random disturbances,
	all of which vary across devices and rounds.
	
	Each service round proceeds through four phases.  In the
	\emph{scheduling} phase~(P1), the DASH Scheduler computes the
	selected set~$\Sset^{(t)}$ and allocations~$\{v_i^*\}$,
	incorporating straggler probabilities
	(Section~\ref{sec:sa_select}).  In the \emph{dispatch} phase~(P2),
	the ES broadcasts $(\mathbf{w}^{(t-1)}, v_i^*)$ to each
	$i \in \Sset^{(t)}$.  In the \emph{collection} phase~(P3), devices
	execute the local pipeline and upload logits while the Async
	Controller enforces deadline~$D^{(t)}$; devices that have not
	finished are either partially collected or timed out
	(Section~\ref{sec:straggler}).  Finally, in the \emph{aggregation}
	phase~(P4), the Aggregation Engine computes quality-weighted logit
	averages over the received contributions and distills them into
	$\mathbf{w}^{(t)}$ (Section~\ref{sec:aggregation}).
	
	% ── 3.2 Device Cost and Quality Model ────────────────────
	\subsection{Device cost and quality model}\label{sec:cost_quality}
	
	Each device~$i$ is characterized by a profile
	$(\rho_i, \theta_i, b_i, a_i)$.  The marginal resource cost for
	uploading one logit vector at round~$t$ is
	\begin{equation}\label{eq:cost}
		b_i^{(t)} \triangleq \omega_i^{(t)} + \mu_i,
	\end{equation}
	where $\omega_i^{(t)} > 0$ is the per-vector communication cost
	(determined by channel conditions and allocated resource blocks under
	OFDMA) and $\mu_i > 0$ is the per-vector computation cost.  The
	activation cost $a_i > 0$ captures the fixed overhead of local model
	training; the total cost for uploading~$v_i$ vectors is thus
	$C_i(v_i) = a_i \cdot \mathbf{1}_{\{v_i > 0\}} + b_i \, v_i$.
	
	The effective quality of device~$i$'s contribution is modeled by a
	saturation function:
	\begin{equation}\label{eq:quality}
		q_i(v_i) = \rho_i \cdot \frac{v_i}{v_i + \theta_i},
	\end{equation}
	where $\theta_i > 0$ is the half-saturation constant (the volume at
	which device~$i$ attains quality $\rho_i / 2$) and
	$\rho_i \in (0, 1]$ is the privacy degradation factor quantifying
	the fraction of useful information retained after local perturbation
	($\rho_i = 1$: no perturbation; $\rho_i \to 0$: strong privacy).
	This Michaelis--Menten form captures diminishing returns
	($\partial q_i / \partial v_i > 0$,
	$\partial^2 q_i / \partial v_i^2 < 0$), bounded capacity
	($q_i \to \rho_i$), and privacy attenuation (the ceiling scales
	with~$\rho_i$).
	
	Under a per-round budget $B > 0$, the ES jointly selects a device
	subset $\Sset \subseteq \Mset$ and determines upload volumes
	$\{v_i\}_{i \in \Sset}$ to maximize aggregate quality:
	\begin{subequations}\label{eq:det_problem}
		\begin{align}
			\max_{\Sset,\,\{v_i\}} \;\;
			& \sum_{i \in \Sset} \rho_i \,\frac{v_i}{v_i + \theta_i}
			\label{eq:det_obj} \\
			\text{s.t.} \;\;
			& \sum_{i \in \Sset}(a_i + b_i \, v_i) \leq B,
			\label{eq:det_budget} \\
			& v_i \in [0, \vmax],\; \forall\, i \in \Sset,
			\;\; \Sset \subseteq \Mset.
			\label{eq:det_bound}
		\end{align}
	\end{subequations}
	Problem~\eqref{eq:det_problem} is a mixed-integer program that is
	NP-hard in general, but admits efficient approximation through the
	structural properties exploited in Section~\ref{sec:scheduling}.
	
	% ── 3.3 Straggler Latency Model ─────────────────────────
	\subsection{Straggler latency model}\label{sec:straggler}
	
	In a synchronous protocol, the ES waits for all selected devices to
	finish before aggregation; the round time is therefore dominated by
	the slowest device.  As~$M$ grows, the probability of at least one
	extreme straggler increases rapidly, making synchronous operation
	increasingly impractical.
	
	We model the round-trip latency of device~$i$ as the sum of three
	components:
	\begin{equation}\label{eq:latency}
		\tau_i = \tau_i^{\mathrm{comp}} + \tau_i^{\mathrm{comm}}
		+ \tau_i^{\mathrm{noise}},
	\end{equation}
	where
	$\tau_i^{\mathrm{comp}} = c_i \cdot |\Dpriv| \cdot E$ is the local
	computation time, with~$c_i$ the per-sample processing rate and~$E$
	the number of local training epochs (deterministic,
	device-dependent);
	$\tau_i^{\mathrm{comm}} = v_i^* \cdot s / r_i$ is the uplink
	transmission time, with~$s$ the payload size per logit vector
	and~$r_i$ the instantaneous channel rate (deterministic given the
	allocation~$v_i^*$); and
	$\tau_i^{\mathrm{noise}} \sim \mathrm{LogNormal}(\mu_n, \sigma_n)$
	captures random disturbances from queueing, OS scheduling, and
	transient interference (stochastic, i.i.d.\ across devices and
	rounds).  The log-normal model for~$\tau_i^{\mathrm{noise}}$ is
	well-established in the networking literature for representing
	heavy-tailed latency variations; the parameter~$\sigma_n$ controls
	the severity of straggler behavior.
	
	At each round~$t$, the ES sets a deadline~$D^{(t)}$
	(Section~\ref{sec:timeout}).  Each selected device~$i$ falls into
	exactly one of three outcome categories:
	
	\begin{definition}[Device Outcome]\label{def:outcome}
		Given deadline~$D^{(t)}$ and device latency~$\tau_i$, device~$i$'s
		outcome is classified as follows.
		\begin{enumerate}
			\item[(C)] Complete ($\tau_i \leq D^{(t)}$): the ES receives all
			$v_i^*$ logit vectors, i.e., $v_i^{\mathrm{recv}} = v_i^*$.
			\item[(P)] Partial
			($\tau_i^{\mathrm{comp}} \leq D^{(t)} < \tau_i$): the device has
			finished computation but its upload is interrupted; the ES receives
			\begin{equation}\label{eq:partial}
				v_i^{\mathrm{recv}} = \Bigl\lfloor v_i^*
				\cdot \frac{D^{(t)} - \tau_i^{\mathrm{comp}}}
				{\tau_i^{\mathrm{comm}}} \Bigr\rfloor
			\end{equation}
			logit vectors.
			\item[(T)] Timeout ($D^{(t)} < \tau_i^{\mathrm{comp}}$): the device
			has not completed local computation; no logits are received
			($v_i^{\mathrm{recv}} = 0$).
		\end{enumerate}
	\end{definition}
	
	The completion probability of device~$i$ under deadline~$D$ is
	\begin{equation}\label{eq:pi}
		\pi_i(D) \triangleq \Pr[\tau_i \leq D]
		= F_{\mathrm{LN}}\!\bigl(D - \tau_i^{\mathrm{comp}}
		- \tau_i^{\mathrm{comm}};\; \mu_n, \sigma_n\bigr),
	\end{equation}
	where~$F_{\mathrm{LN}}$ is the log-normal CDF evaluated at the
	residual slack $D - \tau_i^{\mathrm{comp}} - \tau_i^{\mathrm{comm}}$
	(set to zero if the slack is non-positive).  This quantity plays a
	central role in straggler-aware scheduling
	(Section~\ref{sec:sa_select}).
	
	\begin{definition}[Straggler Rate]\label{def:straggler_rate}
		The per-round straggler rate is the fraction of selected devices that
		fail to deliver complete logits by the deadline:
		\begin{equation}\label{eq:straggler_rate}
			p^{(t)} = \frac{|\{i \in \Sset^{(t)} :
				\tau_i > D^{(t)}\}|}{|\Sset^{(t)}|}.
		\end{equation}
	\end{definition}
	
	% ── 3.4 Asynchronous Protocol ────────────────────────────
	\subsection{Asynchronous protocol}\label{sec:async_protocol}
	
	Algorithm~\ref{alg:async_kaas} formalizes the per-round operation of
	the DASH framework.  Compared to a synchronous counterpart, the
	protocol introduces three modifications: (i)~a deadline~$D^{(t)}$ is
	computed before dispatch (line~2), (ii)~device completion
	probabilities $\pi_i(D^{(t)})$ are incorporated into the scheduling
	decision (line~3), and (iii)~after the deadline expires, only complete
	and (optionally) partial contributions are aggregated
	(lines~8--12).  All logits are computed using the current-round
	global model~$\mathbf{w}^{(t-1)}$, so there is no model
	staleness---a deliberate design choice that keeps the aggregation
	semantics clean and avoids the need for staleness correction weights.
	
	\begin{algorithm}[t]
		\caption{DASH: Asynchronous FD Scheduling Protocol}
		\label{alg:async_kaas}
		\begin{algorithmic}[1]
			\REQUIRE Device pool $\Mset$, budget $B$, rounds $T$,
			reference dataset $\Dref$, timeout policy $\mathcal{P}$
			\ENSURE Global model $\mathbf{w}^{(T)}$
			\STATE Initialize $\mathbf{w}^{(0)}$;\;
			$\mathcal{H} \gets \emptyset$
			\COMMENT{latency history}
			\FOR{$t = 1$ \TO $T$}
			\STATE $D^{(t)} \!\gets\!
			\mathcal{P}.\mathrm{Deadline}(t, \mathcal{H})$
			\STATE $(\Sset^{(t)}\!, \{v_i^*\}) \!\gets\!
			\mathrm{DASH\text{-}Select}(\Mset, B, D^{(t)})$
			\STATE Dispatch $(\mathbf{w}^{(t-1)}, v_i^*)$
			to each $i \in \Sset^{(t)}$
			\STATE Wait until wall-clock time reaches $D^{(t)}$
			\STATE // Collect responses
			\STATE $\Srecv \gets \emptyset$
			\FOR{each $i \in \Sset^{(t)}$}
			\IF{outcome$(i)$ = Complete}
			\STATE $\Srecv \!\gets\! \Srecv \cup \{(i, v_i^*)\}$
			\ELSIF{outcome$(i)$ = Partial
				\AND $\mathcal{P}$ accepts partial}
			\STATE $\Srecv \!\gets\! \Srecv
			\cup \{(i, v_i^{\mathrm{recv}})\}$
			\ENDIF
			\ENDFOR
			\STATE // Aggregate and distill
			\IF{$\Srecv \neq \emptyset$}
			\STATE Compute $\bar{\mathbf{z}}^{(t)}$
			via Eq.~\eqref{eq:agg_weight}
			\STATE Update $\mathbf{w}^{(t)}$
			via Eq.~\eqref{eq:kl}
			\ELSE
			\STATE $\mathbf{w}^{(t)} \!\gets\! \mathbf{w}^{(t-1)}$
			\COMMENT{skip round}
			\ENDIF
			\STATE $\mathcal{H} \!\gets\! \mathcal{H}
			\cup \{\tau_i : i \in \Sset^{(t)}\}$
			\ENDFOR
			\RETURN $\mathbf{w}^{(T)}$
		\end{algorithmic}
	\end{algorithm}
	
	
	% ══════════════════════════════════════════════════════════
	\section{Scheduling and Aggregation}\label{sec:scheduling}
	% ══════════════════════════════════════════════════════════
	
	This section presents the algorithmic components invoked by
	Algorithm~\ref{alg:async_kaas}.  Section~\ref{sec:waterfill} derives
	the optimal volume allocation for a given device set.
	Section~\ref{sec:sa_select} introduces straggler-aware device
	selection, the core theoretical contribution of this paper.
	Section~\ref{sec:timeout} describes three timeout policies, and
	Section~\ref{sec:aggregation} details the quality-weighted
	aggregation.
	
	% ── 4.1 Volume Allocation via Water-Filling ──────────────
	\subsection{Volume allocation via water-filling}\label{sec:waterfill}
	
	For a fixed device subset~$\Sset$ with residual budget
	$\Bres = B - \sum_{i \in \Sset} a_i$, the volume allocation
	sub-problem seeks to maximize aggregate quality:
	\begin{equation}\label{eq:subproblem}
		\max_{\{v_i\}} \;
		\sum_{i \in \Sset} \rho_i \,\frac{v_i}{v_i + \theta_i}
		\;\;\text{s.t.}\;\;
		\sum_{i \in \Sset} b_i v_i \leq \Bres,\;
		v_i \in [0, \vmax].
	\end{equation}
	
	\begin{proposition}[Water-Filling Allocation]\label{prop:wf}
		The optimal solution to~\eqref{eq:subproblem} is
		\begin{equation}\label{eq:wf_sol}
			v_i^* = \min\!\Biggl\{
			\Biggl[\sqrt{\frac{\rho_i\,\theta_i}{\nu^*\, b_i}}
			- \theta_i \Biggr]^{\!+}\!,\; \vmax
			\Biggr\}, \;\; \forall\, i \in \Sset,
		\end{equation}
		where $[x]^+ = \max(x, 0)$ and $\nu^* > 0$ is the unique water
		level satisfying $\sum_{i \in \Sset} b_i v_i^* = \Bres$.
	\end{proposition}
	
	\begin{pf}
		The objective is strictly concave in each~$v_i$
		($\partial^2 q_i / \partial v_i^2 < 0$) and the constraint set is
		convex.  Slater's condition holds whenever $\Bres > 0$.  Writing the
		Lagrangian and setting $\partial \mathcal{L} / \partial v_i = 0$
		yields $\rho_i \theta_i / (v_i + \theta_i)^2 = \nu b_i$, which
		gives~\eqref{eq:wf_sol}.  Uniqueness of~$\nu^*$ follows from the
		intermediate value theorem, since
		$\sum_i b_i v_i^*(\nu)$ is continuous and strictly decreasing
		in~$\nu$.  The water level is computed via bisection in
		$O(|\Sset|\log(1/\delta))$ time for tolerance~$\delta$.
		\qed\end{pf}
	
	The efficiency index
	$\eta_i \triangleq \rho_i / (b_i \theta_i)$ captures the marginal
	quality gain per unit cost at zero allocation.
	From~\eqref{eq:wf_sol}, device~$i$ receives a positive allocation if
	and only if $\eta_i > \nu^*$; devices below this threshold are
	effectively excluded by the water-filling mechanism.
	
	% ── 4.2 Straggler-Aware Device Selection ─────────────────
	\subsection{Straggler-aware device selection}\label{sec:sa_select}
	
	In a synchronous setting, the ES selects devices to maximize the
	deterministic set quality
	$Q(\Sset) = \max_{\{v_i\}} \sum_{i \in \Sset} q_i(v_i)$ subject to
	budget constraint~\eqref{eq:det_budget}, solved by
	Proposition~\ref{prop:wf}.  In the asynchronous setting, however,
	each device~$i$ only delivers its complete contribution with
	probability~$\pi_i(D)$ (Eq.~\ref{eq:pi}).  The ES should therefore
	maximize the expected aggregate quality after straggler realization.
	
	Let $X_i \sim \mathrm{Bernoulli}(\pi_i(D))$ denote the indicator
	that device~$i$ completes on time.  The realized quality is the
	random variable
	$Q^{\mathrm{real}}(\Sset) = \sum_{i \in \Sset} X_i \cdot
	q_i(v_i^*)$.  By linearity of expectation:
	\begin{equation}\label{eq:expected_quality}
		\mathbb{E}\bigl[Q^{\mathrm{real}}(\Sset)\bigr]
		= \sum_{i \in \Sset} \pi_i(D) \cdot q_i(v_i^*)
		\triangleq \tilde{Q}(\Sset).
	\end{equation}
	We call $\tilde{Q}(\Sset)$ the expected quality function, which
	replaces~$Q(\Sset)$ as the optimization objective in the asynchronous
	setting.  Formally, $\tilde{Q}$ is obtained by solving the modified
	allocation problem:
	\begin{subequations}\label{eq:async_problem}
		\begin{align}
			\tilde{Q}(\Sset)
			&\triangleq \max_{\{v_i\}}
			\sum_{i \in \Sset}
			\underbrace{\pi_i(D) \cdot \rho_i}_{\tilde{\rho}_i}
			\,\frac{v_i}{v_i + \theta_i}
			\label{eq:async_obj} \\
			&\;\;\text{s.t.}\;\;
			\textstyle\sum_{i \in \Sset}(a_i + b_i v_i) \leq B,\;
			v_i \in [0, \vmax].
			\label{eq:async_budget}
		\end{align}
	\end{subequations}
	This is structurally identical to~\eqref{eq:det_problem} with the
	effective fidelity $\tilde{\rho}_i = \pi_i(D) \cdot \rho_i$
	replacing~$\rho_i$.  Consequently, Proposition~\ref{prop:wf} applies
	directly with~$\tilde{\rho}_i$ in place of~$\rho_i$, and the
	straggler-aware efficiency index becomes
	$\tilde{\eta}_i = \tilde{\rho}_i / (b_i \theta_i)
	= \pi_i(D) \cdot \eta_i$.
	The preservation of structural properties follows immediately.
	
	\begin{proposition}[Submodularity under Straggler Awareness]
		\label{prop:sub_async}
		$\tilde{Q}: 2^{\Mset} \to \mathbb{R}_{\geq 0}$ is monotone
		non-decreasing and submodular.
	\end{proposition}
	
	\begin{pf}
		For any fixed deadline~$D$, the completion
		probability~$\pi_i(D)$ is a non-negative constant associated with
		each device~$i$.  Problem~\eqref{eq:async_problem} is therefore an
		instance of~\eqref{eq:det_problem} with modified fidelity parameters
		$\tilde{\rho}_i = \pi_i(D) \cdot \rho_i \in (0, \rho_i]$.
		Monotonicity and submodularity of the parametric family
		$Q_{\rho}(\Sset) = \max_{\{v_i\}} \sum_{i \in \Sset}
		\rho_i' v_i / (v_i + \theta_i)$ subject to a knapsack constraint
		hold for any $\rho_i' > 0$: adding a device to a smaller set leaves
		more residual budget, and the concavity of each~$q_i$ implies weakly
		larger marginal gain \citep{boyd_convex}.
		\qed\end{pf}
	
	\begin{theorem}[Approximation Guarantee for $\tilde{Q}$]
		\label{thm:approx_async}
		The greedy algorithm (Algorithm~\ref{alg:dash_select}) returns
		$\Sset^{\mathrm{g}}$ satisfying
		\begin{equation}\label{eq:greedy_bound}
			\tilde{Q}(\Sset^{\mathrm{g}})
			\geq \tfrac{1}{2}(1 - 1/e) \cdot \tilde{Q}^*,
		\end{equation}
		where
		$\tilde{Q}^* = \max_{\Sset \subseteq \Mset} \tilde{Q}(\Sset)$
		is the optimal expected quality.
	\end{theorem}
	
	\begin{pf}
		$\tilde{Q}$ is monotone submodular
		(Proposition~\ref{prop:sub_async}) and the budget
		constraint~\eqref{eq:async_budget} is a knapsack constraint.  The
		result follows from the $\tfrac{1}{2}(1 - 1/e)$ approximation ratio
		for greedy maximization of monotone submodular functions under a
		knapsack constraint \citep{khuller_knapsack}.
		\qed\end{pf}
	
	We now relate the expected quality~$\tilde{Q}^*$ to the
	deterministic optimum $Q^* = \max_{\Sset} Q(\Sset)$, which represents
	the quality achievable if all devices always complete on time.
	
	\begin{theorem}[Expected Quality Bound under Straggler Rate]
		\label{thm:straggler_bound}
		Let $\pi_{\min} = \min_{i \in \Mset} \pi_i(D)$ and define the
		worst-case straggler probability $\bar{p} = 1 - \pi_{\min}$.
		Then the greedy solution satisfies
		\begin{equation}\label{eq:straggler_bound}
			\tilde{Q}(\Sset^{\mathrm{g}})
			\geq (1 - \bar{p}) \cdot
			\tfrac{1}{2}(1 - 1/e) \cdot Q^*.
		\end{equation}
	\end{theorem}
	
	\begin{pf}
		Let~$\Sset^*$ denote the optimal device set for the deterministic
		problem, with water-filling allocation $\{v_i^{\circ}\}$ achieving
		$Q^* = \sum_{i \in \Sset^*} q_i(v_i^{\circ})$.
		
		The allocation $\{v_i^{\circ}\}$ is feasible for
		problem~\eqref{eq:async_problem} with $\Sset = \Sset^*$ (same budget
		constraint).  Therefore:
		\begin{align}
			\tilde{Q}(\Sset^*)
			&\geq \sum_{i \in \Sset^*}
			\pi_i(D) \cdot \rho_i
			\frac{v_i^{\circ}}{v_i^{\circ} + \theta_i}
			\notag \\
			&\geq \pi_{\min}
			\sum_{i \in \Sset^*}
			\rho_i \frac{v_i^{\circ}}{v_i^{\circ} + \theta_i}
			= \pi_{\min} \cdot Q^*.
			\label{eq:step1}
		\end{align}
		The first inequality uses the feasibility of~$\{v_i^{\circ}\}$ (the
		$\tilde{Q}$-optimal water-filling can only yield a higher value); the
		second applies $\pi_i(D) \geq \pi_{\min}$ for all~$i$.
		Since $\tilde{Q}^* \geq \tilde{Q}(\Sset^*)$, we have
		$\tilde{Q}^* \geq \pi_{\min} \cdot Q^*$.
		Combining with Theorem~\ref{thm:approx_async}:
		\begin{align}
			\tilde{Q}(\Sset^{\mathrm{g}})
			&\geq \tfrac{1}{2}(1 - 1/e) \cdot \tilde{Q}^*
			\notag \\
			&\geq (1 - \bar{p}) \cdot
			\tfrac{1}{2}(1 - 1/e) \cdot Q^*.
			\label{eq:step2}
		\end{align}
		\qed\end{pf}
	
	The bound~\eqref{eq:straggler_bound} quantifies the price of
	asynchrony: the greedy expected quality degrades by at most a factor
	of $(1 - \bar{p})$ relative to the deterministic optimum.  This is a
	worst-case bound based on~$\pi_{\min}$; in practice, the
	straggler-aware selection preferentially includes reliable devices
	(high~$\pi_i$), so the effective average completion probability
	across~$\Sset^{\mathrm{g}}$ is substantially higher
	than~$\pi_{\min}$.  The modified efficiency index
	$\tilde{\eta}_i = \pi_i \cdot \rho_i / (b_i \theta_i)$ naturally
	balances reliability against knowledge quality: a device with
	excellent data quality ($\rho_i$ high, $\theta_i$ low) but poor
	connectivity ($\pi_i$ low) is down-ranked relative to a moderately
	informative but reliable peer.
	
	Algorithm~\ref{alg:dash_select} differs from a standard
	deterministic greedy scheduler in exactly one line: the marginal gain
	computation multiplies by~$\pi_i(D)$.  All other steps---ranking,
	water-filling, greedy loop---remain identical.
	
	\begin{algorithm}[t]
		\caption{DASH-Select: Straggler-Aware Scheduling}
		\label{alg:dash_select}
		\begin{algorithmic}[1]
			\REQUIRE Device pool $\Mset$, budget $B$, deadline $D$,
			profiles $\{(\rho_i, \theta_i, a_i, b_i)\}$,
			straggler model $(\mu_n, \sigma_n)$, tolerance $\delta$
			\ENSURE Selected set $\Sset$, allocations $\{v_i^*\}$
			\STATE // Compute straggler-aware efficiency
			\FOR{each $i \in \Mset$}
			\STATE $\pi_i \gets \Pr[\tau_i \leq D]$
			via Eq.~\eqref{eq:pi}
			\STATE $\tilde{\eta}_i \gets
			\pi_i \cdot \rho_i / (b_i \theta_i)$
			\ENDFOR
			\STATE Sort devices by $\tilde{\eta}_i$ (decreasing)
			\STATE // Greedy selection on $\tilde{Q}$
			\STATE $\Sset \gets \emptyset$
			\FOR{each device $i$ in sorted order}
			\IF{$a_i + \sum_{j \in \Sset} a_j > B$}
			\STATE \textbf{continue}
			\ENDIF
			\STATE $\Sset' \gets \Sset \cup \{i\}$
			\STATE $\{v_j\} \!\gets\!
			\mathrm{WaterFill}(\Sset',
			B {-} \textstyle\sum_{j \in \Sset'}\! a_j,
			\{\tilde{\rho}_j\}, \delta)$
			\IF{$\tilde{Q}(\Sset') > \tilde{Q}(\Sset)$}
			\STATE $\Sset \gets \Sset'$
			\ENDIF
			\ENDFOR
			\STATE // Final allocation
			\STATE $\{v_i^*\}_{i \in \Sset} \!\gets\!
			\mathrm{WaterFill}(\Sset,
			B {-} \textstyle\sum_{i \in \Sset}\! a_i,
			\{\tilde{\rho}_i\}, \delta)$
			\RETURN $\Sset,\; \{v_i^*\}$
		\end{algorithmic}
	\end{algorithm}
	
	% ── 4.3 Timeout Policy Design ────────────────────────────
	\subsection{Timeout policy design}\label{sec:timeout}
	
	The deadline~$D^{(t)}$ controls the trade-off between round latency
	(shorter deadline leads to faster rounds) and knowledge utilization
	(longer deadline allows more devices to complete, yielding higher
	per-round quality).  We consider three policies with increasing
	adaptiveness.
	
	The simplest strategy, which we call the fixed-deadline policy, sets
	$D^{(t)} = D_0$ for all rounds~$t$, where $D_0$ is a system-level
	constant chosen by the operator.  This provides predictable per-round
	wall-clock time and is straightforward to implement, but cannot adapt
	to changing network conditions.  Setting $D_0$ too low wastes
	scheduled resources through excessive timeouts, while setting it too
	high approaches synchronous behavior.
	
	The adaptive-deadline policy adjusts the deadline based on observed
	latency statistics from the previous round.  Let
	$\{\tau_i^{(t-1)}\}$ denote the latencies observed in round~$t{-}1$.
	The adaptive policy sets:
	\begin{equation}\label{eq:adaptive}
		D^{(t)} = \alpha_{\mathrm{ema}} \,
		\mathrm{Pctl}_p\!\bigl(\{\tau_i^{(t-1)}\}\bigr)
		+ (1 {-} \alpha_{\mathrm{ema}}) \, D^{(t-1)},
	\end{equation}
	where $\mathrm{Pctl}_p$ is the $p$-th percentile of the empirical
	latency distribution and $\alpha_{\mathrm{ema}} \in (0,1)$ is an
	exponential moving average smoothing coefficient that prevents
	oscillation.  During a warmup phase of~$W$ rounds, the policy falls
	back to $D^{(t)} = D_0$.  The percentile parameter~$p$ directly
	controls the target completion rate: $p = 0.7$ targets approximately
	70\% on-time completion.
	
	The partial-accept policy uses the same adaptive deadline mechanism,
	but additionally accepts partial logit uploads from devices in the
	Partial outcome category (Definition~\ref{def:outcome}).  Partial
	contributions enter the aggregation with reduced quality weight since
	$q_i(v_i^{\mathrm{recv}}) < q_i(v_i^*)$ whenever
	$v_i^{\mathrm{recv}} < v_i^*$.  This policy maximizes knowledge
	utilization by retaining computation that the device has already
	performed, at the cost of slightly noisier aggregation from
	incomplete logit sets.
	
	% ── 4.4 Quality-Weighted Aggregation ─────────────────────
	\subsection{Quality-weighted aggregation}\label{sec:aggregation}
	
	After the deadline expires, the ES aggregates logits from the
	set of responding devices~$\Srecv$, which includes all devices with
	Complete outcomes and, under the partial-accept policy, those with
	Partial outcomes as well.  Each device $i \in \Srecv$ has delivered
	$v_i^{\mathrm{recv}}$ logit vectors (equal to~$v_i^*$ for complete
	devices, or given by~\eqref{eq:partial} for partial devices).  The
	aggregation weight for device~$i$ is
	\begin{equation}\label{eq:agg_weight}
		w_i = \frac{\rho_i \cdot q_i(v_i^{\mathrm{recv}})}
		{\displaystyle\sum_{j \in \Srecv}
			\rho_j \cdot q_j(v_j^{\mathrm{recv}})},
	\end{equation}
	where $q_i(v_i^{\mathrm{recv}}) =
	\rho_i \, v_i^{\mathrm{recv}} / (v_i^{\mathrm{recv}} + \theta_i)$.
	Devices with higher privacy fidelity and more complete uploads
	receive proportionally larger weight, while partial contributions are
	naturally down-weighted through the saturation function.  The
	aggregated logit vector
	$\bar{\mathbf{z}}^{(t)} = \sum_{i \in \Srecv}
	w_i \tilde{\mathbf{z}}_i$
	is then used to update the global model via KL-divergence
	minimization:
	\begin{equation}\label{eq:kl}
		\mathbf{w}^{(t)} \!=\! \mathbf{w}^{(t-1)}
		\!-\! \alpha \nabla_{\mathbf{w}}
		\mathrm{KL}\!\bigl(
		\sigma(\bar{\mathbf{z}}^{(t)} / \tau)
		\,\|\,
		\sigma(\mathbf{w}^{(t-1)}\!(\Dref) / \tau)
		\bigr),
	\end{equation}
	where $\sigma$ is the softmax function, $\tau$ is the distillation
	temperature, and~$\alpha$ is the learning rate.
	
	In the synchronous protocol, $v_i^{\mathrm{recv}} = v_i^*$ for all
	selected devices and $\Srecv = \Sset^{(t)}$, so the aggregation
	weight reduces to $w_i \propto \rho_i \cdot q_i(v_i^*)$, recovering
	the deterministic quality weighting.  In the asynchronous setting,
	partial devices contribute with $v_i^{\mathrm{recv}} < v_i^*$,
	resulting in lower~$q_i$ and hence lower weight---a graceful
	degradation mechanism that requires no additional tuning parameters.
	
	
	% ══════════════════════════════════════════════════════════
	\section{Performance Evaluation}\label{sec:eval}
	% ══════════════════════════════════════════════════════════
	
	% [PLACEHOLDER — 5.1 Setup, 5.2 Main Comparison,
	%  5.3 Straggler Sweep, 5.4 Timeout Policy,
	%  5.5 Scalability, 5.6 EMNIST, 5.7 Privacy]
	%
	% Figures to include:
	%  Fig 2: accuracy vs round (6 methods)
	%  Fig 3: accuracy vs wall-clock time (6 methods) — key figure
	%  Fig 4: accuracy vs straggler ratio
	%  Fig 5: wall-clock vs straggler ratio
	%  Fig 6: Pareto front (policy comparison)
	%  Fig 7: accuracy vs M at fixed wall-clock budget
	%  Fig 8: round time vs M
	%  Fig 9: EMNIST accuracy vs wall-clock
	%  Fig 10: privacy sweep under async
	%
	% Tables:
	%  Table 1: performance summary (6 methods x 3 metrics)
	
	
	% ══════════════════════════════════════════════════════════
	\section{Discussion}\label{sec:discuss}
	% ══════════════════════════════════════════════════════════
	
	% [PLACEHOLDER — 6.1 Sync vs Async guidelines,
	%  6.2 Deadline Tuning, 6.3 Limitations and Future Work]
	
	
	% ══════════════════════════════════════════════════════════
	\section{Conclusion}\label{sec:conclude}
	% ══════════════════════════════════════════════════════════
	
	% [PLACEHOLDER]
	
	
	% ── CRediT ────────────────────────────────────────────────
	\printcredits
	
	% ── Acknowledgment ────────────────────────────────────────
	% AI-Generated Content Disclosure:
	% The Claude AI system (Anthropic) was used to assist with LaTeX
	% formatting and grammatical refinement during manuscript preparation.
	% All technical content, theoretical derivations, experimental design,
	% and scientific conclusions were produced by the authors.
	
	% ── Bibliography ──────────────────────────────────────────
	\bibliographystyle{cas-model2-names}
	
	%% Inline bibliography for standalone compilation.
	%% Replace with \bibliography{refs} once .bib file is ready.
	\begin{thebibliography}{99}
		
		\bibitem[{Boyd and Vandenberghe(2004)}]{boyd_convex}
		Boyd, S., Vandenberghe, L., 2004.
		Convex Optimization.
		Cambridge University Press.
		
		\bibitem[{Khuller et~al.(1999)}]{khuller_knapsack}
		Khuller, S., Moss, A., Naor, J.S., 1999.
		The budgeted maximum coverage problem.
		Information Processing Letters 70~(1), 39--45.
		
		\bibitem[{McMahan et~al.(2017)}]{mcmahan_fl}
		McMahan, B., Moore, E., Ramage, D., Hampson, S., y~Arcas, B.A.,
		2017.
		Communication-efficient learning of deep networks from decentralized
		data.
		In: Proc.\ AISTATS. pp.\ 1273--1282.
		
		\bibitem[{Li and Wang(2019)}]{li_fedmd}
		Li, D., Wang, J., 2019.
		{FedMD}: Heterogeneous federated learning via model distillation.
		arXiv preprint arXiv:1910.03581.
		
		\bibitem[{Hinton et~al.(2015)}]{hinton_kd}
		Hinton, G., Vinyals, O., Dean, J., 2015.
		Distilling the knowledge in a neural network.
		arXiv preprint arXiv:1503.02531.
		
		\bibitem[{Xie et~al.(2019)}]{xie_fedasync}
		Xie, C., Koyejo, S., Gupta, I., 2019.
		Asynchronous federated optimization.
		arXiv preprint arXiv:1903.03934.
		
		\bibitem[{Nguyen et~al.(2022)}]{nguyen_fedbuff}
		Nguyen, J., Malik, K., Zhan, H., Yousefpour, A., Rabbat, M.,
		Malek, M., Huba, D., 2022.
		Federated learning with buffered asynchronous aggregation.
		In: Proc.\ AISTATS. pp.\ 3581--3607.
		
		\bibitem[{Kairouz et~al.(2021)}]{kairouz_fl}
		Kairouz, P., McMahan, H.B., Avent, B., et~al., 2021.
		Advances and open problems in federated learning.
		Foundations and Trends in Machine Learning 14~(1--2), 1--210.
		
		\bibitem[{Gad et~al.(2024)}]{gad_fedskd}
		Gad, G., Fadlullah, Z.M., Fouda, M.M., Ibrahem, M.I., Kato, N.,
		2024.
		Federated learning with selective knowledge distillation over
		bandwidth-constrained wireless networks.
		In: Proc.\ IEEE ICC. pp.\ 3476--3481.
		
		\bibitem[{Balakrishnan et~al.(2022)}]{balakrishnan_diverse}
		Balakrishnan, R., Li, T., Zhou, T., Himayat, N., Smith, V.,
		Bilmes, J., 2022.
		Diverse client selection for federated learning via submodular
		maximization.
		In: Proc.\ ICLR.
		
	\end{thebibliography}
	
\end{document}